\chapter{Methodology \& Tools}
\label{ch:methodology}

\section{Data Collection Pipeline}
\label{sec:data-collection}

\begin{table}[htbp]
\centering
\caption{Data Collection Pipeline Components}
\label{tab:data-collection-pipeline}
\begin{tabular}{@{}ll@{}}
\toprule
\textbf{Component} & \textbf{Description} \\
\midrule
\textbf{Queries} & 79 product recommendation queries $\times$ 3 runs each (237 total runs) \\
\textbf{ChatGPT Data} & \makecell[l]{Full responses with inline citations, additional links, and recommended \\ products through chrome dev tool network packet inspection} \\
\textbf{Bing Data} & Top 200 results \\
\textbf{Gemini Data} & Full \texttt{groundingMetadata} (Chunks vs.\ Supports) + Fan-out Queries \\
\textbf{Google SERP} & \makecell[l]{SerpApi pagination until $\geq$20 Organic results are collected \\ (often $\sim$3 pages), with Video/PAA/Discussions retained \\ as diagnostic buckets} \\
\textbf{Content Fetching} & \makecell[l]{Node.js fetcher + Browser extension for blocked pages \\ (Master Content Library)} \\
\bottomrule
\end{tabular}
\end{table}

\subsection{Selection of the Research Query Set}
\label{sec:query-set-selection}

\begin{itemize}
    \item \textbf{High-Volume Real-World Prompts:} Our dataset consists of \textbf{79 unique product recommendation prompts} (e.g., ``Best AI video translators'', ``Top-rated transcription software''), sourced from two platforms: \textbf{65 prompts from Profound} (real user conversations with ChatGPT) and \textbf{14 prompts from Ahrefs} (high-volume keyword clusters).
    \item \textbf{Methodology for Selection:}
    \begin{itemize}
        \item \textbf{Keyword Clustering:} Using \textbf{Ahrefs} to identify high-intent keyword clusters and derive 14 representative prompts from search volume data.
        \item \textbf{Prompt Volume Analysis:} Leveraging \textbf{Profound's} Prompt Volumes feature to select 65 prompts derived from real user conversations with ChatGPT. Note: Profound states that surfaced prompts may be rewritten for anonymity or clarity, so these are representative of real user intent rather than verbatim transcripts.
        \item \textbf{Deliberate Intent Filtering:} From the broad set of available user prompts, we \textbf{deliberately filtered for high commercial intent}. This ensures the study reflects the specific segment of search where AI synthesis is most active and where the ``Extractive Nature'' of the model is most visible.
        \item \textbf{Domain Expertise:} Queries were focused on the \textbf{AI and Software-as-a-Service (SaaS)} sectors --- particularly \textbf{speech and language technology} (speech-to-text, live transcription, voice translation, text-to-speech) --- a domain where the author has significant professional expertise, allowing for more nuanced qualitative analysis of the ``Signal vs.\ Noise'' in results.
    \end{itemize}
    \item \textbf{Experimental Rigor:} Each of the 79 prompts was executed in \textbf{3 independent runs} (with a 4th run added only in cases of technical failure or RAG non-triggering) to analyze the consistency and stochastic nature of the retrieval process.
\end{itemize}

\subsection{Google SERP Result Types (SerpApi): Organic vs Video vs PAA}
\label{sec:google-serp-types}

SerpApi returns Google results in multiple \textbf{result\_type} buckets (not just ``10 blue links''). This matters because overlap numbers can shift depending on what we count as ``the SERP.''

\begin{itemize}
    \item \textbf{Organic results}: standard web results (our primary control-group baseline).
    \item \textbf{Video results}: often YouTube-heavy; can appear in top positions and inflate ``coverage'' for topics where ChatGPT cites YouTube.
    \item \textbf{Discussions / forums blocks}: SerpApi often surfaces forum-like ``discussions'' sections; we retain them as a separate diagnostic bucket (forums / quota etc.).
\end{itemize}

\textbf{Method rule (comparability)}:
\begin{itemize}
    \item For overlap metrics, we default to \textbf{Organic-only} (and treat Video/PAA as separate diagnostic buckets), unless explicitly stated otherwise.
    \item We keep the non-organic buckets available as a \textbf{discussion point} (e.g., ``Google surfaces YouTube via Video blocks earlier than Bing''), and as a sensitivity analysis (``organic-only vs organic+video'').
    \item \textbf{Pagination rule (how we collected the ``Top-20 Organic'' baseline)}: we paginate SerpApi until we have $\geq$20 Organic results (not necessarily only the first page). In practice this is often the first $\sim$3 pages, alongside non-organic blocks.
\end{itemize}

\subsection{Bing SERP Scraping (Browser Extension)}
\label{sec:bing-scraping}

Because Microsoft retired public Bing Search API access on \textbf{August 11, 2025} (see Section~1.2.5), no stable programmatic API was available for academic SERP collection. We therefore built a \textbf{Chrome browser extension} (\texttt{tools/Bing Results Scraper}) that scrapes Bing's live consumer UI via DOM interaction.

\begin{itemize}
    \item \textbf{Scrape method}: The extension runs as a Manifest V3 content script that simulates human search interaction --- typing queries character-by-character with randomized inter-key delays (10--50\,ms), clicking the search button, and waiting for DOM stabilization (result count checked across 3 consecutive iterations before proceeding).
    \item \textbf{Ad filtering}: Organic results are selected via \texttt{.b\_algo:not(.b\_adTop):not(.b\_adBottom):not([data-apurl])}. A supplementary \texttt{isSponsoredResultElement()} function provides multi-layer ad detection: DOM class markers (\texttt{.adsMvCarousel}, \texttt{.b\_ads1line}, \texttt{.ad\_em}), text-prefix matching (\texttt{/\^{}Sponsored\textbackslash b/i}), CSS \texttt{::before} pseudo-element inspection, and URL-based filtering of known ad-click endpoints (\texttt{/aclick}, \texttt{/clk}, \texttt{/sclk} on \texttt{bing.com}). Unresolved Bing redirect URLs (\texttt{bing.com/ck/}) are also discarded as probable ads.
    \item \textbf{Pagination to Rank 200}: The extension clicks Bing's native pagination button (\texttt{.sb\_pagN}) and waits 3 seconds between pages. Each page yields $\sim$10 organic results; 20 pages $\times$ 10 results = 200 results per query. Positions are numbered strictly sequentially across pages (not reset per page).
    \item \textbf{URL decoding}: Bing wraps organic URLs in redirect parameters. The extension decodes these via URL-parameter extraction, Base64 decoding (with prefix stripping), and hex-byte fallback. Tracking parameters (\texttt{utm\_*}, \texttt{fbclid}, \texttt{gclid}) are stripped to normalize URLs for overlap matching.
    \item \textbf{US proxy}: All Bing scrapes were routed through a \textbf{US-based proxy} to match the retrieval environment used for ChatGPT runs (see Section~2.7.3), ensuring that localized SERP variation does not inflate or deflate overlap measurements.
    \item \textbf{Extracted fields per result}: \texttt{position}, \texttt{title}, \texttt{url}, \texttt{domain}, \texttt{displayUrl}, \texttt{snippet}, \texttt{page\_num}. When content extraction is enabled, the extension also fetches each page and extracts full text (up to 20,000 characters), metadata (\texttt{canonical\_url}, \texttt{published\_date}, \texttt{modified\_date}), schema markup types, and table counts.
\end{itemize}

\section{Methodological Evolution: From Top 30 to ``Deep Hunt'' (Rank 200)}
\label{sec:methodological-evolution}

\begin{itemize}
    \item \textbf{Starting Point (Single-Query, Top 30):} Our study initially used only the \textbf{single fan-out query visible in ChatGPT's UI} (the ``Searching for\ldots'' text) and scraped the \textbf{Top 30 Bing results} for each, yielding a \textbf{40.3\% citation overlap}. At this stage we had no knowledge of the hidden double fan-out mechanism --- we assumed the UI-displayed query was the only query issued.
    \item \textbf{The Network Instrumentation Turning Point:} Once we began capturing raw network payloads (see Section~2.3), we discovered that ChatGPT always dispatches \textbf{two} fan-out queries per search turn, not one. The ``Searching for\ldots'' text in the UI only surfaces one of them. This meant our initial Bing scrapes were matching against an incomplete query set, understating true overlap.
    \item \textbf{The Discovery of ``UI Erasure'':} Even after correcting for fan-out, qualitative review using our custom Data Viewer revealed a significant ``Visibility Gap.'' ChatGPT was citing high-quality, relevant pages that were completely missing from the Top 30 human-facing results.
    \item \textbf{The Pivot to Rank 200:} To test whether these citations were truly ``invisible'' or merely ``buried,'' we expanded our methodology to a \textbf{``Deep Hunt'' (Rank 200)}.
    \item \textbf{Key Finding of the Pivot:} We discovered that Bing often surfaces the exact pages ChatGPT cites, but hides them deep within pagination loops or beyond the ``Page 2 Cliff'' (Rank 11+). This methodological shift allowed us to prove that the difference between Search and GenAI is often a \textbf{UI and Ranking problem}, not just an indexing one.
    \item \textbf{Why Even Rank 200 May Not Be Enough:} Our page-level distribution data (see Section~3.2.1) shows that citation matches do not reach zero at Rank 200 --- the long tail continues. Going deeper (Rank 300+) would likely recover additional matches and further shrink the ``truly invisible'' set, producing a cleaner separation between sources ChatGPT finds via search indices and sources it accesses through other means (parametric memory, direct partnerships, or non-public indices).
\end{itemize}

\subsection{Bing Pagination Instability (The ``Elastic Page 1'' Problem)}
\label{sec:bing-pagination-instability}

A critical methodological challenge arose from Bing's inconsistent UI pagination behavior, which we discovered through repeated scrapes of the same query. This instability is partly why the retirement of the Bing Search API (see Section~1.2.5, Section~2.1.3) is consequential for grounding research: the consumer UI we were forced to scrape is a far noisier baseline than a stable API would have provided.

\begin{itemize}
    \item \textbf{Page 1 Truncation:} Bing's first page of organic results is not a fixed ``Top 10.'' In some scrapes, Page 1 returned as few as \textbf{2 organic results} (e.g., for \texttt{"free AI tools to translate video"}, only \texttt{zeemo.ai} and \texttt{airmore.ai} appeared as organic results on Page 1, while subsequent scrapes of the same query minutes later returned \texttt{maestra.ai}, \texttt{akool.com}, \texttt{veed.io}, \texttt{happyscribe.com}, and \texttt{rask.ai} in the top positions). The rest of the page was filled with ads, Copilot branding, ``People also search for'' blocks, and video carousels.
    \item \textbf{The ``Missing Middle'' Problem:} When Page 1 is truncated, the results at ranks 3--7 do not simply shift to Page 2. Instead, they appear to be \textbf{dropped entirely} from that particular scrape. Page 2 begins with a different, loosely-indexed set of results that often includes irrelevant content (e.g., PDF translators, image-to-video tools, and unrelated AI generators appearing for a video translation query). Our citation-overlap data corroborates this qualitative observation: Page 1 concentrates the bulk of matches, Page 2 shows a conspicuous dip, and Pages 3+ settle into a gradual decay --- the expected long-tail pattern, but at higher match volumes than a tightly indexed SERP would produce. The Page 2 dip in particular stands out as an artifact of loose indexing rather than smooth ranking decay (see Section~3.2.1).
    \item \textbf{Implications for Overlap Measurement:}
    \begin{enumerate}
        \item \textbf{Lost overlap (conservative bias):} If our Bing scrape for a given run captured a truncated Page 1, we may have missed legitimate top-ranking results that ChatGPT's retrieval system almost certainly had access to. This means our reported overlap percentages are likely \textbf{underestimates} --- some citations we classified as ``invisible'' may have been present in Bing's index at high ranks but absent from our specific scrape due to pagination instability.
        \item \textbf{Why Top 200 was necessary:} Because everything after Page 1 is loosely indexed and inconsistent, we needed to scrape deep (Rank 200) not because ChatGPT is truly ``deep hunting'' at those ranks, but to \textbf{compensate for Bing's UI-layer instability} and recover matches that a stable API would have returned in the Top 10--30.
    \end{enumerate}
    \item \textbf{The API vs.\ UI Divergence Hypothesis:} We strongly suspect that ChatGPT does \textbf{not} receive the same noisy, truncated pagination that Bing's web UI serves to human users. Bing likely provides its index to ChatGPT via a backend integration that returns a clean, stable ranked list without ads, carousels, or pagination artifacts --- the exact mechanism is not publicly documented, but the ranking and filtering logic almost certainly differs from what the consumer UI exposes. This means our Bing-UI-scraped baseline is a \textbf{degraded proxy} for what ChatGPT actually sees --- further supporting the interpretation that our overlap figures are conservative lower bounds.
    \item \textbf{What This Implies About ChatGPT's Retrieval:} Given that ChatGPT consistently cites high-quality, relevant sources, one of two things must be true: either ChatGPT receives a \textbf{cleaner, more stable ranking} from Bing's backend API than what the consumer UI exposes, or the model is doing substantial work to \textbf{filter through the same noisy index} and extract signal from noise. Either interpretation underscores the divergence between the human-facing and agent-facing search experiences described in Section~1.2.5.
    \item \textbf{Contrast with Google (SerpApi):} This instability was not observed in our Google SERP data collected via SerpApi, which returned consistent, stable organic rankings across paginated requests. This asymmetry between Bing's UI instability and Google's API stability is itself a noteworthy finding for researchers attempting to replicate grounding studies.
\end{itemize}

\section{Anatomy of a ChatGPT Response (Network-Instrumented)}
\label{sec:chatgpt-anatomy}

\textit{What exactly we can observe about ChatGPT's retrieval + citation pipeline from captured network payloads.}

\subsection{Pipeline Overview (Methodology)}
\label{sec:pipeline-overview}

Our network instrumentation captured raw event-stream payloads from the ChatGPT production UI, revealing a \textbf{three-layer pipeline} (Sonic Classifier $\rightarrow$ Sonicberry Orchestrator $\rightarrow$ Generation Model) rather than a monolithic model. The detailed architecture, layer descriptions, and observed flow are reported as empirical findings in Section~3.1. Here we note only the methodological implication: because each layer emits distinct payload fields, we can separately instrument the search-trigger decision, the fan-out query dispatch, and the citation-generation phase.

\subsection{Source Categories (Naming Conventions)}
\label{sec:source-categories}

We classify every URL that appears in ChatGPT's network response into one of three categories, derived from distinct locations in the payload:

\begin{itemize}
    \item \textbf{Cited}: URLs that appear as inline footnote references in the generated answer. Mechanically, these are URLs linked via \texttt{content\_references} token spans --- the model inserted a citation marker at a specific position in the output text pointing to this source. In the ChatGPT UI, these render as numbered superscript links within the response body.
    \item \textbf{Additional}: URLs that are attached to the response but \textbf{not} referenced inline. These appear in the \texttt{sources\_additional} array in the network payload and render in the ChatGPT UI as a collapsible ``Sources'' section below the main response. The model retrieved and surfaced them to the user, but did not tie them to any specific claim.
    \item \textbf{Unreferenced}: URLs present in the \texttt{search\_result\_groups} payload (the raw search results delivered to the model during generation) that were \textbf{not} promoted to either cited or additional status. These never appear in the user-facing response. We compute them as: all URLs in \texttt{search\_result\_groups} minus those already classified as cited or additional. See Section~3.3.3 for detailed analysis.
\end{itemize}

\textbf{Important:} The \texttt{search\_result\_groups} field is the retrieval pool as it appears in ChatGPT's network stream. We do not know the exact upstream source of these results --- while ChatGPT is known to use Bing, the field itself does not identify the backing search provider, and we cannot rule out internal indices or other retrieval paths.

\subsection{Search Trigger Instrumentation (the ``decision'' fields)}
\label{sec:search-trigger}

We log the system's search-decision artifacts where present (Enterprise streams are richest):
\begin{itemize}
    \item \texttt{sonic\_classification\_result}: the classifier output used to decide search (e.g., \texttt{simple\_search\_prob}, \texttt{complex\_search\_prob}, thresholds).
    \item \texttt{web\_search\_triggered}: whether search actually ran.
    \item \texttt{web\_search\_forced}: whether search was forced (if present).
\end{itemize}

The full Sonic Classifier configuration (classifier variant, 3-class thresholds, Enterprise vs.\ Personal divergence, and comparison with external reverse-engineering by Resoneo) is reported as empirical findings in Section~3.1.1.

\subsection{Fan-Out Queries (``hidden queries'') \& the \texttt{web.run} Agentic Loop}
\label{sec:fan-out-queries}

\begin{itemize}
    \item \textbf{Definition}: multiple search queries issued for one user prompt; these expand/reshape retrieval scope.
    \item \textbf{Why it matters}: fan-out sets control which sources are even eligible to be cited.
    \item \textbf{Observable field (network-derived)}: \texttt{search\_model\_queries.queries} (stored as \texttt{hidden\_queries\_json} in our extracted CSV/DB pipeline).
\end{itemize}

\textbf{The \texttt{web.run} mechanism (empirically discovered):}
\begin{itemize}
    \item The generation model (gpt-5-2) calls a tool named \texttt{web.run} via the Sonicberry orchestrator (\texttt{alpha.sonic\_thinky\_v1\_paid}). Each \texttt{web.run} call \textbf{always dispatches exactly 2 queries in parallel} --- a fixed pair-generation strategy producing one synonym variant and one structural rephrase. Odd query counts (1, 3, 5) are architecturally impossible.
    \item In the vast majority of runs (\textbf{94\%}), a single \texttt{web.run} call is all that fires --- the 2 queries run in parallel and the results are used directly. However, in a small fraction of runs (2.7\%), the generation model evaluates the returned results and calls \texttt{web.run} again, creating an iterative re-search loop tracked by the \texttt{search\_turns\_count} field.
    \item \textbf{How it appears in the raw network stream (multi-turn fan-out signature)}:
    \begin{itemize}
        \item The response arrives as an event stream (patch/append style) containing repeated tool messages (\texttt{role="tool"}, \texttt{name="web.run"}, \texttt{author.metadata.source="sonic\_tool"}).
        \item Each \texttt{web.run} message carries \texttt{search\_model\_queries.queries[]} (always length 2) and \texttt{search\_turns\_count} (incrementing 1, 2, 3\ldots).
        \item Multi-turn calls may be \textbf{chained} (each \texttt{parent\_id} = previous \texttt{web.run}'s \texttt{message\_id}) or \textbf{independent} (different parent).
    \end{itemize}
\end{itemize}

The empirical distribution of search turns, re-search strategies, query drift, and timing evidence are reported in Section~3.1.3.

\subsection{Retrieved Candidate Pool (what the model could have used)}
\label{sec:retrieved-candidates}

We capture the retrieved candidates and their metadata:
\begin{itemize}
    \item \texttt{search\_result\_groups}: grouped search entries containing \texttt{url}, \texttt{title}, \texttt{snippet}, and \texttt{ref\_id} (turn/ref index keys).
    \item \textbf{Interpretation}: this is the model's ``shortlist menu'' prior to selection.
\end{itemize}

\subsection{Citation Tokens \& Claim Mapping (how we attach sources to text)}
\label{sec:citation-tokens}

Two key payload elements enable claim$\rightarrow$source reconstruction:
\begin{itemize}
    \item \texttt{content\_references}: token spans / citation tokens with \texttt{start\_idx} / \texttt{end\_idx} (where citations appear in the output stream).
    \item \texttt{response\_text}: the generated text (often streamed via patch/append events).
\end{itemize}

From these, we build:
\begin{itemize}
    \item \textbf{Claim blocks}: the text segments preceding each citation token (our \texttt{claim\_text} extraction).
    \item \textbf{Mapped sources}: resolved URLs/titles/snippets by matching \texttt{ref\_id} keys to \texttt{search\_result\_groups}.
\end{itemize}

\subsection{What the Network Does \textit{Not} Reveal (critical limitation)}
\label{sec:network-limitations}

\begin{itemize}
    \item It does \textbf{not} include the exact on-page snippets/chunks that were fetched/inserted into the model context.
    \item Therefore: our ChatGPT-side ``grounding budget'' analyses are \textbf{output-side proxies} (claim-attributed text), not true input-side snippet budgets.
\end{itemize}

\subsection{What Happens Next (how this section feeds the thesis)}
\label{sec:anatomy-next-steps}

This ``anatomy'' motivates the next analytic layers:
\begin{itemize}
    \item \textbf{Overlap \& visibility}: compare cited/additional/unreferenced pools against Bing Top 30 + Deep Hunt and Google SERP controls.
    \item \textbf{Selection bias}: compare Content DNA of Cited vs Additional vs Unreferenced.
    \item \textbf{Stochasticity}: quantify fan-out drift across runs and resulting citation churn.
\end{itemize}

\subsection{Network Parameter Glossary (ChatGPT, Network-Instrumented)}
\label{sec:network-glossary}

\textit{A practical glossary of the recurring fields we use when decoding the raw event stream. This is intentionally limited to parameters we directly observed in captured payloads.}

\subsubsection{Identifiers / linkage}
\label{sec:glossary-identifiers}

\begin{description}
    \item[\texttt{conversation\_id}] Conversation session identifier (useful for grouping, not analysis).
    \item[\texttt{message.id}] Unique message identifier in the stream.
    \item[\texttt{parent\_id}] Links a message to its parent; helps follow the chain of events.
    \item[\texttt{request\_id}] Correlates events belonging to the same request; often shared across tool calls and patches for that run.
    \item[\texttt{turn\_exchange\_id} / \texttt{turn\_trace\_id}] Internal tracing IDs for a single turn; useful for debugging continuity across events.
\end{description}

\subsubsection{Search trigger \& decision}
\label{sec:glossary-search-trigger}

\begin{description}
    \item[\texttt{sonic\_classification\_result}] The search trigger classifier output (probabilities + thresholds).
    \item[\texttt{web\_search\_triggered} / \texttt{web\_search\_forced}] Whether search ran / was forced (when present).
    \item[\texttt{message\_marker}] Markers like \texttt{search\_start} that indicate when retrieval begins.
\end{description}

\subsubsection{Fan-out queries \& search turns}
\label{sec:glossary-fan-out}

\begin{description}
    \item[\texttt{metadata.search\_model\_queries.queries[]}] The model-generated fan-out query list for a search turn (our ``hidden queries'').
    \item[\texttt{search\_turns\_count}] Which search turn we are in (1, 2, 3\ldots); multi-turn fan-out appears as repeated query batches with increasing counts.
    \item[\texttt{search\_tool\_call\_count}] Run-level count of search tool invocations; typically correlates with \texttt{max(search\_turns\_count)} (single-turn: 1; multi-turn: 2+).
    \item[\texttt{search\_source} / \texttt{client\_reported\_search\_source}] Indicates the origin mode (often \texttt{composer\_auto}); useful for auditing when behavior changes.
\end{description}

\subsubsection{Retrieved candidates (what came back from search)}
\label{sec:glossary-retrieved}

\begin{description}
    \item[\texttt{metadata.search\_result\_groups}] The grouped retrieval pool (URLs, titles, snippets).
    \item[\texttt{ref\_id}] The key (turn/ref indices) we use to map candidates to citations (\texttt{turn\_index}, \texttt{ref\_type}, \texttt{ref\_index}).
    \item[\texttt{debug\_sonic\_thread\_id}] Internal identifier for the search thread (debugging/trace only).
\end{description}

\subsubsection{Citation / attribution artifacts}
\label{sec:glossary-citation}

\begin{description}
    \item[\texttt{content\_references}] Citation token spans inserted into the streamed response; indices point to the citation token location, not the claim boundary.
    \item[\texttt{safe\_urls} and \texttt{url\_moderation} events] Safety/allow-list decisions for URLs that show up in citations and link cards.
\end{description}

\subsubsection{Completion / response framing}
\label{sec:glossary-completion}

\begin{description}
    \item[\texttt{finish\_details}] How generation ended (stop reason / tokens).
    \item[\texttt{model\_slug} / \texttt{default\_model\_slug}] The model variant used for the run.
\end{description}

\section{Anatomy of a Gemini Response (API-Instrumented)}
\label{sec:gemini-anatomy}

\textit{How the Gemini Vertex AI API exposes grounding metadata compared to ChatGPT's hidden network stream.}

\subsection{The \texttt{groundingMetadata} Schema}
\label{sec:grounding-metadata}

Unlike ChatGPT, where we must ``scrape'' the network stream, Gemini provides structured grounding data in the API response:
\begin{itemize}
    \item \textbf{\texttt{webSearchQueries}}: The explicit fan-out query set generated by the model.
    \item \textbf{\texttt{groundingChunks}}: The raw snippets of text retrieved from the web (the ``injected context'').
    \item \textbf{\texttt{groundingSupports}}: The precise mapping between specific segments of the response and the \texttt{groundingChunks} (the ``paper trail'').
\end{itemize}

\subsection{Key Differences in Instrumentation}
\label{sec:gemini-instrumentation-differences}

\begin{itemize}
    \item \textbf{Transparency}: Gemini is ``Grounding-First''---it exposes the raw chunks it read, whereas ChatGPT only exposes the final URL and a snippet.
    \item \textbf{Segment-Level Attribution}: Gemini attributes every sentence/segment to a specific chunk index, allowing for a much higher resolution of fidelity analysis.
    \item \textbf{Vertex Redirects}: Gemini's \texttt{groundingChunks} do not expose raw URLs --- instead they contain internal redirect URLs (e.g., \texttt{vertexaisearch.cloud.google.com/\ldots}) that must be followed to discover the actual destination domain. We built a dedicated resolution script (\texttt{scripts/resolve\_grounding\_urls.mjs}) that collects all unique Vertex redirect URLs from the raw response files, follows each redirect via HTTP, and produces a mapping from Vertex URL $\rightarrow$ resolved destination URL. Some redirects failed to resolve programmatically (e.g., due to timeouts or anti-bot blocks), so unresolved URLs were exported and resolved manually via a browser-based content fetcher. Without this resolution step, no domain-level or overlap analysis would be possible on the Gemini data.
    \item \textbf{Thinking budget}: All Gemini runs in this study used the \textbf{minimum thinking budget} available via the API. Higher thinking budgets produce more fan-out queries --- the model generates an initial broad query, discovers brands/products from the results, then issues targeted follow-up queries for those specific entities (e.g., \texttt{"best free AI video translation tools 2025 2026"} $\rightarrow$ \texttt{"HeyGen free trial video translation"} $\rightarrow$ \texttt{"ElevenLabs video translator free tier limitations"}). This iterative refinement pattern resembles GPT's multi-turn re-search (see Section~3.1.3), but is driven by the thinking budget rather than an explicit agentic tool-call loop. Our use of minimum thinking keeps the fan-out count bounded and comparable across runs, but means the study captures the model's baseline retrieval behavior rather than its maximum-depth grounding capability.
    \item \textbf{API vs.\ Consumer UI gap}: Our Gemini data comes entirely from the Vertex AI API, which provides structured grounding metadata but does not carry IP/locale signals. Inspecting the \textbf{Gemini consumer UI} (\texttt{gemini.google.com}) via network packet analysis --- mirroring our ChatGPT instrumentation approach --- could reveal additional signals not exposed in the API, such as locale-driven fan-out rewriting, internal ranking or filtering stages before the \texttt{groundingMetadata} is constructed, or differences in fan-out strategy between the consumer product and the API. This remains a potential avenue for future work.
\end{itemize}

\section{Content DNA Enrichment (LLM-as-a-Labeler)}
\label{sec:content-dna}

\textit{How we transformed raw URLs into structured data for selection bias analysis.}

\subsection{The Labeling Pipeline}
\label{sec:labeling-pipeline}

To quantify selection effects (what gets cited vs.\ what was available), we needed to move beyond URLs and domains. We developed an automated enrichment pipeline using \textbf{GPT-5-mini} and \textbf{Gemini Flash 2.5} as structured labelers:
\begin{enumerate}
    \item \textbf{Content Extraction}: Raw HTML was fetched and converted to clean markdown/text.
    \item \textbf{Schema-Driven Labeling}: The LLM was prompted to evaluate each page against a strict 20-field schema, including:
    \begin{itemize}
        \item \textbf{Page Type}: \texttt{listicle}, \texttt{product\_page}, \texttt{documentation}, \texttt{forum\_ugc}, etc.
        \item \textbf{Content Format}: \texttt{best\_of\_list}, \texttt{landing\_page}, \texttt{comparison\_matrix}, etc.
        \item \textbf{Structural Features}: \texttt{has\_tables}, \texttt{has\_numbered\_lists}, \texttt{has\_pros\_cons}.
        \item \textbf{Qualitative Scores}: \texttt{promotional\_intensity\_score}, \texttt{readability\_score}.
        \item \textbf{Tone}: \texttt{promotional}, \texttt{neutral\_informational}, \texttt{salesy}, \texttt{opinionated}.
    \end{itemize}
    \item \textbf{Validation}: A subset of labels was manually audited to ensure the LLM labeler correctly distinguished between vendor-owned landing pages and independent editorial listicles.
\end{enumerate}

\textbf{Note on tone and promotional intensity:} While tone (\texttt{promotional} / \texttt{neutral\_info} / \texttt{salesy} / \texttt{opinionated}) and \texttt{promotional\_intensity\_score} were captured as part of the enrichment schema, they showed no discriminative power in downstream analysis. Tone distribution was uniformly $\sim$70\% promotional across all systems, all citation categories (cited, additional, ignored), and both deployment tiers --- a natural consequence of the product-recommendation prompt set. Because these fields do not differentiate cited from non-cited sources or vary across conditions, they are omitted from the findings tables in Section~3.4--3.5.

\subsection{Enrichment Logic \& Static Overrides}
\label{sec:enrichment-logic}

\textit{To ensure efficiency and accuracy, the labeling pipeline uses a hybrid approach of LLM-labeling and static rules for high-volume, well-known domains.}

\begin{itemize}
    \item \textbf{LLM-Labeling (GPT-5-mini)}: Used for general web pages, blogs, and niche product sites to determine structural features (\texttt{has\_tables}, \texttt{has\_pros\_cons}, etc.).
    \item \textbf{Static Domain Overrides (Skipped Enrichment)}: Known platforms with consistent structural patterns were assigned ``pre-made'' labels to save quota and ensure consistency:
    \begin{itemize}
        \item \textbf{\texttt{reddit.com}}: Automatically labeled as \texttt{type=forum\_ugc}, \texttt{content\_format=discussion\_thread}, \texttt{tone=opinionated}.
        \item \textbf{\texttt{en.wikipedia.org}}: Automatically labeled as \texttt{type=reference}, \texttt{content\_format=encyclopedic}, \texttt{tone=neutral\_informational}.
        \item \textbf{\texttt{arxiv.org}}: Labeled as \texttt{type=reference}, \texttt{content\_format=academic\_paper}.
        \item \textbf{App Stores (\texttt{apps.apple.com}, \texttt{play.google.com})}: Labeled as \texttt{type=app\_store\_listing}, \texttt{primary\_intent=transactional}.
    \end{itemize}
    \item \textbf{The ``Invisible'' Domain Strategy}: Many of the top ``invisible'' domains (like \texttt{reddit.com} and \texttt{wikipedia.org}) were handled via these static overrides because their structure is fixed and does not require per-page LLM analysis.
\end{itemize}

\textbf{Why we skipped these domains from full enrichment:} The skipDomains list (\texttt{wikipedia.org}, \texttt{reddit.com}, \texttt{arxiv.org}, \texttt{github.com}, \texttt{youtube.com}, app stores, social media platforms, and first-party support/docs sites) was applied consistently across all enrichment scripts (\texttt{enrich\_gemini\_grounding.mjs}, \texttt{enrich\_control\_gpt.mjs}, \texttt{enrich\_db\_list\_with\_gemini.js}, \texttt{enrich\_db\_list\_with\_vertex.js}). These domains have fixed, well-understood structures that do not benefit from LLM-based content analysis --- a Wikipedia article is always \texttt{type=reference}, a Reddit thread is always \texttt{type=forum\_ugc}. Instead of wasting API quota on these, we assigned heuristic labels via \texttt{label\_skipped\_urls.mjs} (see static overrides above) with conservative defaults (e.g., \texttt{promotional\_intensity\_score=0}, \texttt{spamminess\_score=0}).

\textbf{Consequence for drift analysis:} Because these skipped domains receive zeroed-out structural scores (no tables, no numbered lists, no pros/cons), including them in feature-level drift comparisons would create artificial signal. The selection drift analysis in Section~3.5 therefore focuses primarily on \textbf{product pages and listicles} --- the two dominant types where full LLM enrichment was performed and where structural features are meaningfully variable. This scoping avoids comparing LLM-enriched pages against heuristically-labeled ones, which would conflate labeling method differences with actual content differences.

\subsection{Labeler Fidelity (cross-model agreement)}
\label{sec:labeler-fidelity}

To test whether Content DNA is model-dependent, we ran a direct agreement audit between \textbf{Gemini 2.5 Flash} and \textbf{GPT-5-mini} on \textbf{2,660 overlapping URLs} (\texttt{docs/inter\_model\_fidelity\_report.md}).

\begin{itemize}
    \item \textbf{Structural fields (agreement)}: \texttt{has\_pros\_cons} \textbf{94.6\%}, \texttt{has\_sources\_or\_citations} \textbf{93.8\%}, \texttt{has\_clear\_authorship} \textbf{93.3\%}, \texttt{has\_tables} \textbf{92.9\%}, \texttt{has\_numbered\_lists} \textbf{91.1\%}
    \item \textbf{Categorical fields (agreement)}: \texttt{content\_format} \textbf{92.1\%}, \texttt{type} \textbf{87.7\%}, \texttt{tone} \textbf{77.7\%}
    \item \textbf{Score consistency (correlation)}: \texttt{freshness\_cue\_strength} \textbf{$r = 0.841$} (strong trend agreement, different absolute thresholds)
\end{itemize}

\subsection{Enrichment Coverage \& Label Distribution (Study URL Universe)}
\label{sec:enrichment-coverage}

To make downstream analyses defensible, we first measured how much of the URL universe was successfully enriched.

\begin{itemize}
    \item \textbf{Enriched URLs:} 11,925 / 11,929 (\textbf{$\sim$100\%})
    \item \textbf{Unlabeled URLs:} 4
\end{itemize}

\subsection{Research Questions for Selection Analysis}
\label{sec:selection-research-questions}

\textit{These research-design questions guided the enrichment-based analyses in Part 3 (Section~3.4--3.5).}

\begin{enumerate}
    \item \textbf{Why were Additional links not cited inline?}
    \begin{itemize}
        \item Compare structural DNA: \texttt{has\_tables}, \texttt{has\_numbered\_lists}, \texttt{heading\_density}
        \item Compare \texttt{tone}: Are Additional links more \texttt{promotional} or \texttt{salesy}?
        \item Compare \texttt{type}: Are Additional links more \texttt{product\_page} vs.\ \texttt{listicle}?
    \end{itemize}

    \item \textbf{Cross-Run Citation:}
    \begin{itemize}
        \item Were Additional links from Run 1 cited inline in Run 2/3/4?
        \item This shows consistency vs.\ randomness in ChatGPT's citation selection
    \end{itemize}

    \item \textbf{Page 1 Ignored Links:}
    \begin{itemize}
        \item Links in \textbf{Bing Page 1} (variable-size SERP page; see Section~3.2.1) that ChatGPT did NOT cite
        \item Compare their DNA to cited links
        \item Hypothesis: Ignored links are more \texttt{salesy}, lower structural complexity
    \end{itemize}
\end{enumerate}

\section{Citation Mapping \& Claim-Level Attribution}
\label{sec:citation-mapping}

\textit{How we precisely map ChatGPT's written claims to their retrieved sources. This pipeline was applied specifically to \textbf{listicle-cited claims}, where the research question is sharpest: when a model cites a listicle containing 10+ products, does it actually extract information from that page, or does it hallucinate from parametric knowledge and merely attach the listicle as a plausible-looking source? Listicles are uniquely suited to this test because they contain a discrete, verifiable roster of items we can check against the model's output. We further restricted the analysis to \textbf{solo-cited claims} --- claim blocks attributed to exactly one URL --- because when multiple sources back a single claim (e.g., a listicle + a product page), it becomes impossible to determine which source the model actually drew from. This solo-citation filter was applied across both GPT and Gemini listicle-cited claims.}

\subsection{The ``Claim-to-Link'' Forensic Pipeline}
\label{sec:claim-to-link}

\begin{itemize}
    \item \textbf{The Challenge:} ChatGPT's final response text replaces internal citation tokens with generic \texttt{[URL]} tags. To understand \textit{why} a link was cited, we must reconstruct the link between the \textbf{written claim} and the \textbf{retrieved source}. Unlike Gemini --- which natively provides segment-level attribution via \texttt{groundingSupports} (see Section~2.4.2) --- ChatGPT exposes no such mapping; we had to build one.
    \item \textbf{The Solution:} We developed a forensic mapping pipeline that reconstructs ChatGPT's claim-to-source attribution at a resolution comparable to Gemini's native \texttt{groundingSupports}:
    \begin{enumerate}
        \item \textbf{Token Alignment:} Raw citation tokens (e.g., \texttt{citeturn0search9}, \texttt{citeturn0search24}) are extracted from the network event stream along with their \texttt{start\_idx} / \texttt{end\_idx} character positions in the streamed response text.
        \item \textbf{Block-Level Extraction:} The script identifies the \textbf{Full Claim Block} --- the descriptive text between consecutive citation tokens. This captures the complete product description or factual statement ChatGPT attributed to that source, not just a keyword. For example, a single claim block might read: \textit{``InqScribe -- desktop transcription software with a perpetual license (no subscription; runs locally)''} mapped to \texttt{citeturn0search24} $\rightarrow$ \texttt{inqscribe.com}.
        \item \textbf{Metadata Enrichment:} Each claim block is joined to its retrieved ``Ground Truth'' (the snippet, title, and URL from the \texttt{search\_result\_groups}) by resolving \texttt{ref\_id} keys (\texttt{turn\_index}, \texttt{ref\_type}, \texttt{ref\_index}).
        \item \textbf{Output:} The pipeline produces a structured \texttt{mapping.json} per run, where each entry contains \texttt{claim\_text}, \texttt{token}, \texttt{target\_url}, \texttt{start\_idx}, and \texttt{end\_idx} --- enabling the same claim-level fidelity and multi-source analyses that Gemini's API provides natively.
    \end{enumerate}
    \item \textbf{Cross-model parity:} This pipeline is what makes direct GPT-vs-Gemini comparison possible at the claim level. Without it, ChatGPT attribution would be limited to run-level URL lists with no way to determine which claim maps to which source.
\end{itemize}

\subsection{Multi-Chip Reconstruction (Synthesis Aggression)}
\label{sec:multi-chip}

\begin{itemize}
    \item \textbf{Defining Multi-Chips:} We observed cases where ChatGPT groups multiple sources under a single citation (e.g., ``Vibe Voice+1'').
    \item \textbf{Forensic Discovery:} Our mapping revealed that these correspond to concatenated tokens (e.g., \texttt{turn0search8} + \texttt{search15}).
    \item \textbf{Research Value:} This allows us to measure \textbf{Synthesis Aggression}---how ChatGPT merges facts from multiple distinct search results into a single cohesive claim.
\end{itemize}

\subsection{Definitions for Multi-Source \& Mixed Citation Analysis}
\label{sec:multi-source-definitions}

\begin{itemize}
    \item \textbf{Multi-cited claim/segment:} a claim/segment with \textbf{2+ distinct cited URLs}. Computed by \texttt{scripts/analysis/multi\_source\_claim\_support.py}.
    \item \textbf{Mixed listicle+product citation:} among multi-cited claims, a case where the cited URL set contains $\geq$1 \texttt{type=listicle} and $\geq$1 \texttt{type=product\_page}.
    \item \textbf{Type-mix buckets:} each multi-cited claim is partitioned into exactly one bucket: \textit{mixed listicle+product}, \textit{listicle-only}, \textit{product-only}, or \textit{neither} (other page types / unknown).
    \item \textbf{Quantitative results} for these metrics are reported in Section~3.6.1.
\end{itemize}

\section{Localization \& Retrieval Environment}
\label{sec:localization}

\textit{How geographical context affects the comparison between Search and GenAI.}

\subsection{Implicit vs.\ Explicit Localization (Definitions \& Instrumentation)}
\label{sec:localization-definitions}

\begin{itemize}
    \item \textbf{Prompt Language Distribution}: Our dataset consists of \textbf{72 English prompts} and \textbf{7 foreign-language prompts} (1 French, 1 Chinese, 2 Turkish, 1 German, 1 Italian, 1 Spanish). All 7 foreign-language prompts originated from \textbf{Ahrefs} keyword clusters; the \textbf{Profound}-sourced prompts were exclusively English.
    \item \textbf{Explicit Localization:} When the query targets an inherently local category --- even without any geographic qualifier in the prompt. A query like ``best bakery'' or ``best pizza'' is implicitly local, and the system will inject a location into the fan-out queries automatically (e.g., rewriting to ``best bakery near Munich Germany''). This primarily applies to local-commerce queries and is not a factor in our SaaS/software dataset, where products are globally available.
    \item \textbf{Implicit Localization:} When the query is general (e.g., ``Best laptop''), but the search engine uses the user's IP, browser language, and search history to localize results.
    \item \textbf{Observed instrumentation signal (Fan-Out Queries):} In our logged \textbf{fan-out query sets} (Gemini \texttt{groundingMetadata.webSearchQueries[]} --- officially called ``grounding support'' queries, functionally equivalent to fan-outs; ChatGPT network-derived hidden queries), implicit localization often manifests as \textbf{one of the fan-out queries being rewritten into the prompt's local language}, which then steers retrieval toward localized sources.
    \item \textbf{Concrete example (implicit localization via fan-out):} When issuing an English prompt from \textbf{Munich, Germany}, we observed a two-query fan-out where one query remained English while the other was rewritten into German:
    \begin{itemize}
        \item Q1: \texttt{"free website or program to translate video and add subtitles"}
        \item Q2: \texttt{"kostenlos video \"ubersetzen und Untertitel automatisch hinzuf\"ugen \ldots"}
    \end{itemize}

    This resulted in one of the two fan-out queries retrieving German-language sources despite an English user prompt.

    \item \textbf{Concrete example (explicit localization via fan-out):} A location-free prompt like \texttt{"what is the best bakery"} can trigger fan-out queries that inject a specific place (e.g., \texttt{"best bakery near Munich Germany"}), effectively converting an implicit prompt into an explicitly localized retrieval task.

    \item \textbf{Observed GPT fan-out localization patterns (62 runs with non-English fan-out queries):}
    \begin{enumerate}
        \item \textbf{Case 1 --- Foreign-language prompts (30 runs):} When the user prompt is in a non-English language (French P019, Chinese P032, Turkish P038, German P055, Turkish P060), GPT emits fan-out queries in both the prompt's language and English. This is expected behavior --- the system mirrors the prompt language while also anchoring to English for broader coverage. Gemini exhibits the same pattern: verified across Chinese (P032: 1 English + 3 Chinese), French (P019: 1 English + 3 French), German (P055: 2 English + 2 German), and Spanish (P062: 1 English + 3 Spanish) prompts.
        \item \textbf{Case 2 --- Context-triggered anomalies (9 runs):} English prompts that \textbf{mention a specific language or locale} in their content trigger non-English fan-out queries. For example, ``best software for Parsian/Farsi audio-to-text transcription'' (P058) produces Farsi fan-out queries; ``Chrome extension for live Chinese to English translation'' (P047) produces Chinese fan-out queries. The non-English query is contextually motivated by the prompt's subject matter, not by IP or browser locale.
        \item \textbf{Case 3 --- Untriggered anomalies (23 runs):} Purely English prompts with \textbf{no language or locale context} produce non-English fan-out queries (Chinese, Japanese, Spanish, French, etc.). For example, ``Is there a real time audio to text translation?'' (P007) or ``Can I add a live translation app to calls?'' (P048) generate Chinese and Japanese fan-out queries. These appear to be driven by IP/locale signals or unexplained model behavior, as nothing in the prompt suggests a non-English retrieval path.
    \end{enumerate}

    \textit{Note: Cases 2 and 3 are GPT-only observations. Gemini was accessed via API without IP-based locale signals, so these patterns could not be tested for Gemini (see Section~2.4.2).}

    \item \textbf{Findings} (occurrence rates, ``English Anchor'' effect, freshness steering stats) are reported in Section~3.1.
\end{itemize}

\subsection{Freshness Steering (Methodology)}
\label{sec:freshness-steering}

\begin{itemize}
    \item \textbf{What we measure:} Whether fan-out queries contain explicit year signals (e.g., ``2025'', ``2026'') and at which query index they appear.
    \item \textbf{Why it matters:} Explicit year injection steers retrieval toward time-stamped listicles, which changes the composition of the grounding pool.
    \item \textbf{Findings} (Gemini vs.\ GPT freshness rates) are reported in Section~3.1.
\end{itemize}

\subsection{The Proxy Requirement (US-Centric Baseline)}
\label{sec:proxy-requirement}

To ensure a fair ``apples-to-apples'' comparison, we standardized our retrieval environment using a \textbf{US-based Proxy}.

\textbf{Why US Proxy?}
\begin{enumerate}
    \item \textbf{Baseline Consistency:} ChatGPT's primary training data and search behaviors are heavily weighted toward US-English web content.
    \item \textbf{Avoiding ``Regional Noise'':} Prevents Bing from surfacing local retailers or regional blogs that ChatGPT would never see, which would artificially lower the overlap percentage.
    \item \textbf{Global Tech Standard:} Most product recommendations in the ``AI/Software'' category (our primary focus) are global in nature, making the US SERP the most relevant ``Ground Truth.''
\end{enumerate}

\subsection{Ethical and Legal Constraints in Retrieval Instrumentation}
\label{sec:ethical-legal}

\begin{itemize}
    \item \textbf{The ``Data Access'' Bottleneck}: A significant challenge in RAG research is the increasing difficulty of accessing ``raw'' search indices.
    \item \textbf{Google vs.\ SerpApi (Dec 2025)}: On December 19, 2025, Google filed a lawsuit against \textbf{SerpApi}, alleging ``unlawful scraping'' and circumvention of security measures\footnote{\url{https://blog.google/innovation-and-ai/technology/safety-security/serpapi-lawsuit/}}.
    \item \textbf{The ``Customer List'' Paradox}: Interestingly, the SerpApi homepage has historically listed major AI players like \textbf{Perplexity} and \textbf{OpenAI} (the latter was subsequently removed) as customers\footnote{\url{https://serpapi.com/}}. This suggests a complex ecosystem where the very companies building RAG systems may rely on third-party scrapers to bridge the ``Visibility Gap'' between their models and the live web.
    \item \textbf{Impact on Methodology}: This legal pressure has led to technical restrictions in the SEO/GEO tool ecosystem, such as the removal of high-volume parameters (e.g., \texttt{num=100}).
    \item \textbf{Research Justification}: These constraints further justify our \textbf{Deep Hunt (Rank 200)} methodology. As traditional scraping becomes more restricted, the ``Visibility Gap'' between what an LLM can see (via direct API access) and what a researcher can see (via public search UIs) will likely widen, making the LLM a primary---and increasingly exclusive---gateway to the deep web.
\end{itemize}

\section{Analysis Applications}
\label{sec:analysis-applications}

We built two separate interactive tools for qualitative inspection and aggregate analysis, one per model:

\subsection{ChatGPT Analysis Viewer (\texttt{scripts/utility/data\_viewer.py})}
\label{sec:chatgpt-viewer}

\begin{itemize}
    \item Per-query/run comparison of ChatGPT responses vs.\ Bing SERP results (Top 200)
    \item Highlights cited, additional, and invisible links with overlap status
    \item Dashboard with aggregate statistics (Overlap \%, Invisible Domains, etc.)
    \item Enabled manual spot-checking that revealed the ``Pagination Loop'' and ``Page 2 Cliff'' problems in Bing
\end{itemize}

\subsection{Gemini Grounding Analyzer (\texttt{tools/GeminiVizApp})}
\label{sec:gemini-viewer}

\begin{itemize}
    \item Per-query/run comparison of Gemini responses vs.\ Google SERP results
    \item Visualizes the grounding pipeline by comparing \texttt{groundingChunks} (retrieved candidates) vs.\ \texttt{groundingSupports} (final segment-level attributions) vs.\ Google SERP (control group)
    \item Surfaces fan-out query overlap distribution per query index (Q1--Q7)
    \item Aggregate grounding behavior dashboard (SERP overlap rate, per-query citation counts)
\end{itemize}

\subsection{Why two apps}
\label{sec:why-two-apps}

The two models expose fundamentally different data structures --- ChatGPT requires reconstructed citation mapping from network tokens (see Section~2.6), while Gemini provides native \texttt{groundingMetadata} with chunk-level and segment-level attribution. A single unified viewer would have obscured these structural differences rather than surfacing them.
